\chapter{Support Vector Machines e metodi kernel}
\label{cap:svm}
\def\svmwA{0.62802}\def\svmwB{1.10490}\def\svmb{-3.95143}

\def\svmswA{1.70187}\def\svmswB{0.58319}\def\svmsb{-4.20225}


\begin{inbreve}
Le Support Vector Machines (SVM) realizzano il secondo approccio alla Structural Risk Minimization. Fissato
l'errore di training, minimizzano automaticamente la VC-confidence: tra tutti gli iperpiani separatori scelgono
quello a \emph{margine massimo}. Il capitolo procede in quattro passi. Prima la SVM lineare \emph{hard margin}, con
la risoluzione tramite moltiplicatori di Lagrange. Poi il \emph{soft margin} per dati non separabili. Poi i
\emph{kernel} per confini non lineari. Infine le SVM per la regressione. Si chiude con una visione critica dell'uso
pratico e con i metodi kernel in generale.
\end{inbreve}

\begin{attenzione}[Notazione]
In questo capitolo si scrive $b$ invece di $w_0$ per il bias. Gli esempi sono indicizzati con $i$ e $j$, da 1 a
$l$. Le parti segnate con $^*$ (i dettagli della risoluzione con i moltiplicatori di Lagrange) sono utili per
completezza, ma non indispensabili per usare il modello.
\end{attenzione}

\section{SVM lineare: il margine massimo}

Una SVM è una \textbf{macchina lineare} (come la LTU) che \textbf{massimizza il margine di separazione}. Così
realizza una \textbf{Structural Risk Minimization}. Inizialmente (\termine{hard margin}) si assume che il problema
sia \textbf{linearmente separabile} e senza errori nei dati.

Dato il training set $T = \{(\x_i,d_i)\}_{i=1}^{l}$ con $d_i\in\{-1,+1\}$, si cerca un iperpiano $\w\T\x+b=0$ con
\[
\w\T\x_i + b\ge0\ \text{ se } d_i=+1,\qquad \w\T\x_i+b<0\ \text{ se } d_i=-1 .
\]
$g(\x) = \w\T\x+b$ è la \termine{funzione discriminante} e $h(\x) = \sign(g(\x))$ l'ipotesi.

\begin{definizione}[Margine di separazione]
Si consideri un iperpiano alla stessa distanza dal più vicino esempio positivo e dal più vicino negativo. Il
\termine{margine di separazione} $\rho$ è il doppio della distanza tra l'iperpiano e il punto più vicino: una
``zona di sicurezza'' attorno al confine. L'\termine{iperpiano ottimo} $\w_o\T\x+b_o=0$ è quello che
\textbf{massimizza il margine}.
\end{definizione}

Non tutti gli iperpiani separatori sono uguali: al variare dell'iperpiano varia il margine. Si vedrà che
$\rho = 2/\norm{\w}$, quindi
\[
\text{massimizzare }\rho\iff\text{minimizzare }\norm{\w}.
\]

\subsection{Rappresentazione canonica e support vector}

Moltiplicando $\w$ e $b$ per una costante l'iperpiano non cambia: è la \textbf{libertà di scala}. Si possono
quindi riscalare $\w$ e $b$ in modo che i punti più vicini all'iperpiano soddisfino
$\abs{g(\x_i)} = \abs{\w\T\x_i+b} = 1$. Allora
\[
\w\T\x_i+b\ge1\ \text{ se } d_i=+1,\qquad \w\T\x_i+b\le-1\ \text{ se } d_i=-1,
\]
cioè, in forma compatta,
\begin{equation}\label{eq:vincoli-svm}
\boxed{\;d_i\,(\w\T\x_i+b)\ge1\qquad\forall i=1,\dots,l\;}
\end{equation}

\begin{definizione}[Support vector]
Un \termine{support vector} $\x^{(s)}$ soddisfa il vincolo \textbf{con l'uguaglianza}:
$d^{(s)}(\w\T\x^{(s)}+b) = 1$. I support vector sono i punti più vicini all'iperpiano, sul bordo del margine.
\end{definizione}

\begin{figure}[ht]
\centering
\begin{tikzpicture}
\begin{axis}[scatola, width=7.4cm, height=6.6cm, xmin=-0.2, xmax=4.6, ymin=-0.2, ymax=4.4, xtick=\empty, ytick=\empty,
  xlabel={$x_1$}, ylabel={$x_2$}, domain=-0.2:4.6, samples=2]
\addplot[thick, inchiostro] {(-\svmb - \svmwA*x)/\svmwB} node[pos=0.06, above right, font=\scriptsize] {$\w\T\x+b=0$};
\addplot[thin, inchiostro, densely dashed] {(1-\svmb - \svmwA*x)/\svmwB} node[pos=0.25, above right, font=\scriptsize] {$+1$};
\addplot[thin, inchiostro, densely dashed] {(-1-\svmb - \svmwA*x)/\svmwB} node[pos=0.3, below left, font=\scriptsize] {$-1$};
\addplot[only marks, mark=*, mark size=1.7pt, blu] table[x=x, y=y, restrict expr to domain={\thisrow{c}}{1:1}] {dati/svmh_punti.dat};
\addplot[only marks, mark=o, mark size=1.9pt, rosso, thick] table[x=x, y=y, restrict expr to domain={\thisrow{c}}{-1:-1}] {dati/svmh_punti.dat};
\addplot[only marks, mark=o, mark size=4pt, inchiostro] table {dati/svmh_sv.dat};
\draw[->, very thick, inchiostro] (axis cs:3.05,1.7) -- (axis cs:{3.05+0.6*\svmwA},{1.7+0.6*\svmwB}) node[right, font=\small] {$\w$};
\end{axis}
\end{tikzpicture}
\caption{L'iperpiano a margine massimo (calcolato su questi dati). Le rette tratteggiate $\w\T\x+b=\pm1$ delimitano
il margine; i punti cerchiati sono i support vector.}
\label{fig:svm-margine}
\end{figure}

\subsection{Calcolo del margine}

\begin{proprieta}{Distanza dall'iperpiano e margine}
La distanza (con segno) di un punto $\x$ dall'iperpiano ottimo è $r = g(\x)/\norm{\w_o}$, e il margine vale
\[
\rho = \frac{2}{\norm{\w_o}} .
\]
\end{proprieta}
\begin{proof}
Poiché $\w_o$ è ortogonale all'iperpiano (\rif{sec:classificazione-lineare}), ogni punto si scrive come
\[
\x = \x_p + r\,\frac{\w_o}{\norm{\w_o}},
\]
dove $\x_p$ è la proiezione di $\x$ sull'iperpiano e $r$ la distanza (\figrif{fig:distanza}). Valutando la
funzione discriminante:
\[
g(\x) = \w_o\T\x_p + b_o + r\,\w_o\T\frac{\w_o}{\norm{\w_o}} = \underbrace{g(\x_p)}_{=0} + r\,\frac{\norm{\w_o}^2}{\norm{\w_o}} = r\,\norm{\w_o},
\]
perché $\x_p$ sta sull'iperpiano. Quindi $r = g(\x)/\norm{\w_o}$. Per un support vector positivo $g(\x^{(s)})=1$:
la sua distanza dall'iperpiano è $1/\norm{\w_o}$. Essa è metà del margine, $\rho/2 = 1/\norm{\w_o}$, da cui
$\rho = 2/\norm{\w_o}$.
\end{proof}

\begin{figure}[ht]
\centering
\begin{tikzpicture}[scale=1.1]
\draw[->, grigio] (-0.3,0) -- (4.2,0) node[right, black, font=\small] {$x_1$};
\draw[->, grigio] (0,-0.3) -- (0,3.4) node[above, black, font=\small] {$x_2$};
\draw[thick] (0.4,3.2) -- (3.6,0.2) node[below, font=\footnotesize] {$\w_o\T\x+b_o=0$};
\coordinate (xp) at (1.9,1.79);
\coordinate (x) at ($(xp)+(0.94*1.25,1.0*1.25)$);
\draw[->, very thick] (xp) -- ($(xp)+(0.47,0.5)$) node[left, font=\small] {$\w_o$};
\draw[densely dashed] (xp) -- (x) node[midway, right, font=\small] {$r$};
\fill (xp) circle (1.6pt) node[below left, font=\small] {$\x_p$};
\fill (x) circle (1.6pt) node[right, font=\small] {$\x$};
\end{tikzpicture}
\caption{La distanza $r$ tra un punto $\x$ e l'iperpiano ottimo: $\x = \x_p + r\,\w_o/\norm{\w_o}$.}
\label{fig:distanza}
\end{figure}

\subsection{Il problema di ottimizzazione}

\begin{definizione}[Problema primale --- hard margin]
Dato $T = \{(\x_i,d_i)\}_{i=1}^l$, trovare $\w$ e $b$ che minimizzano
\[
\Psi(\w) = \frac12\,\w\T\w \qquad(\text{cioè minimizzare }\norm{\w})
\]
soggetto ai vincoli (zero errori di classificazione) $d_i(\w\T\x_i+b)\ge1$ per ogni $i=1,\dots,l$.
\end{definizione}

La funzione obiettivo è \textbf{quadratica e convessa} in $\w$, e i vincoli sono \textbf{lineari} in $\w$. È quindi
un problema di \termine{programmazione quadratica} (QP), che si considera risolto da pacchetti esistenti.
Interessa capire come è formulato, la relazione tra support vector e moltiplicatori, e la forma della soluzione.
Risolverlo in questa forma scala con la dimensione dell'input $n$.

\subsection{Soluzione con i moltiplicatori di Lagrange$^*$}

Si costruisce la \textbf{lagrangiana}
\begin{equation}\label{eq:lagrangiana}
J(\w,b,\balpha) = \frac12\,\w\T\w - \sum_{i=1}^{l}\alpha_i\big(d_i(\w\T\x_i+b)-1\big),
\end{equation}
con un \textbf{moltiplicatore di Lagrange} $\alpha_i\ge0$ per ogni vincolo. $J$ va minimizzata rispetto a $\w$ e
$b$ e massimizzata rispetto ad $\balpha$: la soluzione è un \textbf{punto di sella} di $J$.

\subsubsection{Condizioni di ottimalità}
Derivando $J$ e uguagliando a zero:
\begin{align}
\frac{\partial J}{\partial\w} &= \w - \sum_{i=1}^l\alpha_id_i\x_i = \vv0 &&\Longrightarrow\quad \w = \sum_{i=1}^{l}\alpha_i\,d_i\,\x_i , \label{eq:w-alfa}\\
\frac{\partial J}{\partial b} &= -\sum_{i=1}^l\alpha_id_i = 0 &&\Longrightarrow\quad \sum_{i=1}^{l}\alpha_i\,d_i = 0 . \label{eq:somma-alfa}
\end{align}

\subsubsection{Condizioni di Kuhn-Tucker}
Nel punto di sella vale, per ogni $i$,
\[
\alpha_i\big(d_i(\w\T\x_i+b)-1\big) = 0 .
\]
Quindi: se $\alpha_i>0$ allora $d_i(\w\T\x_i+b) = 1$, cioè $\x_i$ è un \textbf{support vector}. Se $\x_i$ non è un
support vector allora $\alpha_i=0$.

\begin{idea}[L'iperpiano dipende solo dai support vector]
Da \eqref{eq:w-alfa}, $\w_o = \sum_{i=1}^{N_s}\alpha_{o,i}\,d_i\,\x_i$, dove la somma si restringe agli $N_s$
support vector. Tutti gli altri punti si potrebbero rimuovere senza cambiare la soluzione.
\end{idea}

\subsubsection{Il problema duale}
Si sostituiscono le condizioni di ottimalità nella lagrangiana. Sviluppando \eqref{eq:lagrangiana}:
\[
J = \frac12\,\w\T\w - \sum_i\alpha_id_i\,\w\T\x_i - b\sum_i\alpha_id_i + \sum_i\alpha_i .
\]
\begin{enumerate}
  \item Per \eqref{eq:somma-alfa} il terzo termine è nullo.
  \item Per \eqref{eq:w-alfa}, $\w\T\w = \sum_i\sum_j\alpha_i\alpha_jd_id_j\,\x_i\T\x_j$. Anche il secondo termine
        vale $\sum_i\alpha_id_i\big(\sum_j\alpha_jd_j\x_j\big)\T\x_i = \sum_i\sum_j\alpha_i\alpha_jd_id_j\,\x_i\T\x_j$.
  \item Quindi $J = \frac12S - S + \sum_i\alpha_i$, con $S = \sum_{i,j}\alpha_i\alpha_jd_id_j\,\x_i\T\x_j$, e resta una
        funzione dei soli moltiplicatori.
\end{enumerate}

\begin{definizione}[Problema duale --- hard margin]
Trovare $\{\alpha_i\}_{i=1}^l$ che massimizzano
\[
Q(\balpha) = \sum_{i=1}^{l}\alpha_i - \frac12\sum_{i=1}^{l}\sum_{j=1}^{l}\alpha_i\alpha_j\,d_id_j\,\x_i\T\x_j
\]
soggetto a $\alpha_i\ge0$ per ogni $i$ e $\sum_{i=1}^l\alpha_id_i = 0$.
\end{definizione}

Gli $\alpha$ si trovano con la programmazione quadratica o con approcci più efficienti (ad esempio \textbf{SMO},
\ingl{Sequential Minimal Optimization}). Il duale scala con il \textbf{numero di esempi $l$}, meno con la dimensione
dell'input. Si noti che i dati compaiono \textbf{solo tramite prodotti scalari} $\x_i\T\x_j$: è la chiave dei kernel.

\begin{algoritmo}[Ricavare iperpiano e bias]
\begin{enumerate}
  \item Risolvere il duale, ottenendo $\{\alpha_{o,i}\}$.
  \item Calcolare $\w_o = \sum_{i=1}^l\alpha_{o,i}\,d_i\,\x_i$.
  \item Calcolare $b_o = d^{(s)} - \w_o\T\x^{(s)}$ per un support vector $\x^{(s)}$, perché
        $d^{(s)}(\w_o\T\x^{(s)}+b_o) = 1$ e $d^{(s)}=\pm1$. Per un support vector positivo è $b_o = 1-\w_o\T\x^{(s)}$,
        cioè $b_o = 1-\sum_i\alpha_{o,i}d_i\,\x_i\T\x^{(s)}$.
\end{enumerate}
\end{algoritmo}

Non serve conoscere esplicitamente $\w_o$: dato un pattern $\x$ si calcola
\[
g(\x) = \sum_{i=1}^{l}\alpha_{o,i}\,d_i\,\underbrace{\x_i\T\x}_{\text{prodotto scalare}} + b_o
\]
(la somma si restringe ai support vector) e si classifica con il segno di $g(\x)$.

\subsection{Perché il margine massimo generalizza}
\label{sec:svm-srm}

Su un problema separabile l'errore di training è \textbf{fissato} (a zero). Minimizzare $\norm{\w}$ equivale allora
a \textbf{minimizzare la VC-dimension}, e quindi la VC-confidence nel bound $R\le\Remp+\eps(\VC,l,\delta)$
(\rif{sec:vc-bound}).

\begin{teorema}[Vapnik]
Sia $D$ il diametro della più piccola sfera che contiene i dati $\x_1,\dots,\x_l$. Per la classe degli iperpiani
separatori con margine $\rho$ la VC-dimension è limitata da
\[
\VC \le \min\left(\left\lceil\frac{D^2}{\rho^2}\right\rceil,\ n\right) + 1 ,
\]
dove $n$ è la dimensione dell'input.
\end{teorema}

Sia $R_s = D/2$ il raggio della sfera. Con $\rho = 2/\norm{\w}$ si ha $D^2/\rho^2 = R_s^2\,\norm{\w}^2$. Restringendo
la ricerca agli iperpiani a margine massimo, la VC-dimension può quindi essere \textbf{minore di $n+1$}, quella
degli iperpiani generici. \uline{Più il margine è grande, più la VC-dimension è piccola}, indipendentemente dalla
dimensione dello spazio. L'iperpiano a margine massimo fornisce quindi
\begin{itemize}
  \item una \textbf{soluzione unica} (a differenza del Perceptron) con \textbf{zero errori} (a differenza dell'LMS);
  \item una \textbf{SRM automatica}, che minimizza la VC-confidence dentro l'ottimizzazione, \textbf{senza
        iperparametri} nel caso separabile;
  \item la risoluzione con la programmazione quadratica vincolata invece della discesa del gradiente, con una forma
        duale che mette in luce i support vector e i prodotti scalari;
  \item una soluzione concentrata su esempi ``selezionati'', i support vector. È un pregio, perché non dipende dai
        punti lontani dal confine. È un difetto, perché i punti sul confine potrebbero essere i più rumorosi.
\end{itemize}

\section{Soft margin: dati non separabili}

Se almeno un punto viola la condizione $d_i(\w\T\x_i+b)\ge1$ si passa al \termine{soft margin}. Ammettere punti
dentro il margine permette anche un \textbf{margine più ampio}: tolleranza al rumore e modello più semplice.

Si introducono le \termine{variabili slack} $\xi_i\ge0$ e i vincoli diventano
\[
d_i(\w\T\x_i+b)\ge1-\xi_i\qquad\forall i .
\]
\begin{itemize}
  \item $\xi_i=0$: il punto è fuori dal margine o sul bordo;
  \item $0<\xi_i<1$: il punto è dentro il margine ma classificato correttamente;
  \item $\xi_i=1$: il punto sta esattamente sull'iperpiano ($g(\x_i)=0$);
  \item $\xi_i>1$: il punto è classificato male.
\end{itemize}
Un support vector soddisfa ora $d_i(\w\T\x_i+b) = 1-\xi_i$. Il teorema di Vapnik \textbf{non vale più}: è
dimostrato per iperpiani senza alcun punto dentro il margine, e il soft margin toglie proprio questa ipotesi.

\begin{figure}[ht]
\centering
\begin{tikzpicture}
\begin{axis}[scatola, width=7.4cm, height=6.6cm, xmin=-0.6, xmax=4.8, ymin=-0.4, ymax=4.6, xtick=\empty, ytick=\empty,
  xlabel={$x_1$}, ylabel={$x_2$}, domain=-0.6:4.8, samples=2]
\addplot[thick, inchiostro] {(-\svmsb - \svmswA*x)/\svmswB};
\addplot[thin, inchiostro, densely dashed] {(1-\svmsb - \svmswA*x)/\svmswB};
\addplot[thin, inchiostro, densely dashed] {(-1-\svmsb - \svmswA*x)/\svmswB};
\addplot[only marks, mark=*, mark size=1.7pt, blu] table[x=x, y=y, restrict expr to domain={\thisrow{c}}{1:1}] {dati/svms_punti.dat};
\addplot[only marks, mark=o, mark size=1.9pt, rosso, thick] table[x=x, y=y, restrict expr to domain={\thisrow{c}}{-1:-1}] {dati/svms_punti.dat};
\addplot[quiver={u=\thisrow{c}*\thisrow{xi}*\svmswA/(\svmswA^2+\svmswB^2), v=\thisrow{c}*\thisrow{xi}*\svmswB/(\svmswA^2+\svmswB^2)},
  -{Stealth[length=1.5mm]}, arancio, thick, restrict expr to domain={\thisrow{xi}}{0.02:100}] table[x=x, y=y] {dati/svms_punti.dat};
\end{axis}
\end{tikzpicture}
\caption{Soft margin (calcolato su questi dati con $C=1$): le frecce arancioni partono dai punti con $\xi_i>0$ e
arrivano al bordo del margine della loro classe; la loro lunghezza è $\xi_i/\norm{\w}$.}
\label{fig:soft-margin}
\end{figure}

\begin{definizione}[Problema primale --- soft margin]
Trovare $\w$ e $b$ che minimizzano
\[
\Psi(\w,\bxi) = \frac12\,\w\T\w + C\sum_{i=1}^{l}\xi_i
\]
soggetto a $d_i(\w\T\x_i+b)\ge1-\xi_i$ e $\xi_i\ge0$ per ogni $i$.
\end{definizione}

$C$ è un \textbf{iperparametro di regolarizzazione}: si perde la SRM completamente automatica. Regola il
compromesso tra minimizzazione del rischio empirico e del termine di capacità.
\begin{itemize}
  \item \textbf{$C$ basso}: molti errori di training ammessi, margine ampio, possibile \textbf{underfitting}.
  \item \textbf{$C$ alto}: nessun errore ammesso, margine più piccolo, possibile \textbf{overfitting}. Per
        $C\to\infty$ si ritrova l'hard margin.
\end{itemize}

\begin{definizione}[Problema duale --- soft margin$^*$]
Trovare $\{\alpha_i\}$ che massimizzano
\[
Q(\balpha) = \sum_{i=1}^{l}\alpha_i - \frac12\sum_{i=1}^{l}\sum_{j=1}^{l}\alpha_i\alpha_j\,d_id_j\,\x_i\T\x_j
\]
soggetto a $\sum_i\alpha_id_i=0$ e $0\le\alpha_i\le C$ per ogni $i$.
\end{definizione}

È \textbf{identico} al caso hard margin, tranne il vincolo superiore $\alpha_i\le C$. Le condizioni di Kuhn-Tucker
diventano
\[
\alpha_i\big(d_i(\w\T\x_i+b)+\xi_i-1\big) = 0,\qquad \mu_i\,\xi_i = 0,
\]
dove i $\mu_i$ sono i moltiplicatori dei vincoli $\xi_i\ge0$. Ne segue che se $0<\alpha_i<C$ allora $\xi_i=0$, e il
punto è \textbf{sul bordo} del margine. Se $\alpha_i = C$ allora $\xi_i\ge0$, e il punto è \textbf{dentro} il
margine (o dal lato sbagliato). La soluzione è $\w_o = \sum_i\alpha_{o,i}d_i\x_i$. Il bias è
$b_o = d_j - \sum_i\alpha_{o,i}d_i\,\x_i\T\x_j$ per un esempio $j$ con $0<\alpha_j<C$, oppure la media su tutti
questi, per stabilità numerica. I moltiplicatori non nulli corrispondono ai support vector. L'uso è come prima:
$h(\x) = \sign\big(\sum_i\alpha_{o,i}d_i\,\x_i\T\x + b_o\big)$.

\section{Kernel: SVM non lineari}
\label{sec:kernel}

\subsection{Mappare in uno spazio delle feature}

L'idea è mappare i dati dallo spazio degli input a uno \textbf{spazio delle feature} ad alta dimensione, dove
diventano linearmente separabili (si ricordi la LBE, \rif{sec:lbe}).

\begin{esempio}[Da $\R^2$ a $\R^3$]
Si usi $\Phi:\R^2\to\R^3$, $(x_1,x_2)\T\mapsto(x_1^2,\ \sqrt2\,x_1x_2,\ x_2^2)\T$, cioè i monomi di secondo grado.
Due classi separate da un'ellisse nello spazio di input diventano separabili da un \textbf{piano} nello spazio delle
feature (\figrif{fig:phi}). L'equazione dell'ellisse $x_1^2/a^2 + x_2^2/c^2 = 1$ diventa infatti
$z_1/a^2 + z_3/c^2 = 1$, lineare nelle nuove coordinate.
\end{esempio}

\begin{figure}[ht]
\centering
\begin{tikzpicture}
\begin{axis}[scatola, width=5.4cm, height=5.4cm, xmin=-1.7, xmax=1.7, ymin=-1.7, ymax=1.7, xtick=\empty, ytick=\empty,
  xlabel={$x_1$}, ylabel={$x_2$}, title={\small spazio di input}]
\addplot[only marks, mark=*, mark size=1.3pt, rosso] table[x=x, y=y, restrict expr to domain={\thisrow{c}}{1:1}] {dati/phi_2d.dat};
\addplot[only marks, mark=x, mark size=1.8pt, blu] table[x=x, y=y, restrict expr to domain={\thisrow{c}}{0:0}] {dati/phi_2d.dat};
\addplot[thick, inchiostro, domain=0:360, samples=80] ({0.96*cos(x)}, {0.96*0.7*sin(x)});
\end{axis}
\end{tikzpicture}\hspace{0.3cm}
\begin{tikzpicture}
\begin{axis}[width=6.8cm, height=6cm, view={120}{18}, xlabel={$z_1$}, ylabel={$z_2$}, zlabel={$z_3$}, ticks=none,
  title={\small spazio delle feature}, label style={font=\small}, xmin=0, xmax=2.6, ymin=-3.6, ymax=3.6, zmin=0, zmax=2.6]
\addplot3[surf, opacity=0.25, fill=grigio, draw=grigio, domain=0:0.92, y domain=-3.6:3.6, samples=2, samples y=2] (x, y, {(0.92-x)*0.49});
\addplot3[only marks, mark=*, mark size=1.3pt, rosso] table[x=z1, y=z2, z=z3, restrict expr to domain={\thisrow{c}}{1:1}] {dati/phi_3d.dat};
\addplot3[only marks, mark=x, mark size=1.8pt, blu] table[x=z1, y=z2, z=z3, restrict expr to domain={\thisrow{c}}{0:0}] {dati/phi_3d.dat};
\end{axis}
\end{tikzpicture}
\caption{Con la mappa $\Phi(\x) = (x_1^2,\sqrt2x_1x_2,x_2^2)$ le due classi, separate da un'ellisse nello spazio
di input, diventano separabili da un piano (grigio) nello spazio delle feature.}
\label{fig:phi}
\end{figure}

La mappa dell'esempio è costruita apposta per il problema. Trovare la trasformazione migliore \textbf{non è
automatico} (a meno di conoscenza a priori), e spesso si usano trasformazioni generali. \uline{La scelta di $\Phi$ è
il punto critico delle SVM.} Il procedimento ha due passi: (1) mappare i pattern in uno spazio delle feature ad
alta dimensione; (2) cercare l'iperpiano ottimo in quello spazio. Il passo (1) è giustificato dal teorema
seguente.

\begin{teorema}[Cover]
Si proietti in modo non lineare un problema complesso di classificazione di pattern in uno spazio ad alta
dimensione. Lì il problema ha più probabilità di essere linearmente separabile che in uno spazio a bassa
dimensione, purché lo spazio non sia densamente popolato.
\end{teorema}

Un esempio: con una mappa deterministica, $l$ esempi possono essere portati sui vertici di un simplesso di
dimensione $l-1$, il ``triangolo'' generalizzato. Lì \textbf{ogni} partizione binaria degli esempi è linearmente
separabile.

Spazi delle feature molto grandi possono però essere \textbf{computazionalmente infattibili}. Possono anche portare
a \textbf{overfitting} se non si controlla la complessità. L'approccio kernel gestisce lo spazio delle feature
\textbf{implicitamente} e in un contesto regolarizzato: la complessità è regolata dal \textbf{margine}, non dalla
dimensione dello spazio. Con $\Phi:\R^{n}\to\R^{n_\Phi}$ il training set diventa $\{(\Phi(\x_i),d_i)\}$ e
l'iperpiano $\w\T\Phi(\x)+b=0$. Includendo il bias ($w_0=b$, $\phi_0(\x)=1$), $\w\T\Phi(\x)=0$ è un'espansione
lineare in basi in cui la dimensione dello spazio può essere grandissima. Con \eqref{eq:w-alfa} nello spazio delle
feature si ha $\w = \sum_i\alpha_id_i\Phi(\x_i)$, e l'iperpiano diventa
\[
\sum_{i=1}^{l}\alpha_i\,d_i\,\underbrace{\Phi\T(\x_i)\,\Phi(\x)}_{\text{prodotto scalare}} = 0 .
\]
Il prodotto scalare è ora nello spazio delle feature, dove calcolare $\Phi(\x)$ potrebbe essere intrattabile.

\subsection{Il kernel trick}

Sotto certe condizioni \textbf{non serve calcolare $\Phi(\x)$}, e nemmeno conoscere lo spazio delle feature. Basta
una funzione che calcoli direttamente il prodotto scalare in quello spazio.

\begin{definizione}[Kernel]
Un \termine{kernel} (\ingl{inner product kernel}) è una funzione
\[
k:\R^{n}\times\R^{n}\to\R,\qquad k(\x_i,\x) = \Phi\T(\x_i)\,\Phi(\x),
\]
\textbf{simmetrica}: $k(\x_i,\x) = k(\x,\x_i)$.
\end{definizione}

\begin{esempio}[Il kernel dell'esempio]
Con $\Phi(\x) = (x_1^2,\sqrt2x_1x_2,x_2^2)\T$, per $\x = (x_1,x_2)\T$ e $\y = (y_1,y_2)\T$:
\[
\Phi\T(\x)\Phi(\y) = x_1^2y_1^2 + 2x_1x_2y_1y_2 + x_2^2y_2^2 = (x_1y_1+x_2y_2)^2 = (\x\T\y)^2 = k(\x,\y).
\]
Il prodotto scalare nello spazio delle feature (3D) si calcola \textbf{senza passare per la mappa}, direttamente
dai pattern di input (2D).
\end{esempio}

I prodotti scalari tra le immagini degli esempi di training formano la \termine{matrice kernel} (o \termine{matrice di
Gram}) $l\times l$, $K = \{k(\x_i,\x_j)\}_{i,j=1}^l$, simmetrica. Non ogni funzione calcola un prodotto scalare in
qualche spazio. Per il \termine{teorema di Mercer} ciò vale per i kernel che producono matrici kernel
\textbf{semidefinite positive}, cioè con autovalori non negativi. Se $k_1$ e $k_2$ sono kernel, lo sono anche
$k_1+k_2$, $c\,k_1$ (con $c>0$) e $k_1k_2$.

\subsection{SVM con kernel}

Il primale nello spazio delle feature è il soft margin con $\Phi(\x_i)$ al posto di $\x_i$. Ma $\w$ vive nello
spazio delle feature, e il problema può essere intrattabile. Il \textbf{duale} invece resta trattabile:

\begin{definizione}[Problema duale nello spazio delle feature]
Trovare $\{\alpha_i\}$ che massimizzano
\[
Q(\balpha) = \sum_{i=1}^{l}\alpha_i - \frac12\sum_{i,j=1}^{l}\alpha_i\alpha_j\,d_id_j\,\underbrace{k(\x_i,\x_j)}_{K_{ij}}
\]
soggetto a $\sum_i\alpha_id_i = 0$ e $0\le\alpha_i\le C$.
\end{definizione}

Il kernel incorpora la mappa. La dipendenza da $l$ (invece che dalla dimensione dello spazio) diventa un
\textbf{vantaggio}, perché lo spazio delle feature può avere dimensione altissima o infinita.

\begin{algoritmo}[SVM con kernel]
\begin{enumerate}
  \item Dato il training set $T=\{(\x_i,d_i)\}$, scegliere il parametro $C$ e la funzione kernel $k$.
  \item Calcolare la matrice kernel $K$.
  \item Risolvere il problema di ottimizzazione (programmazione quadratica), trovando $\{\alpha_i\}$.
  \item Calcolare il bias $b$ dai moltiplicatori e dalla matrice kernel.
  \item Per un nuovo pattern $\x$ calcolare $\w\T\Phi(\x) + b = \sum_i\alpha_id_i\,k(\x,\x_i) + b$, senza mai
        calcolare $\w$, e classificare con il segno.
\end{enumerate}
\end{algoritmo}

La soluzione è \textbf{sparsa} negli $\alpha$: i moltiplicatori dei non-support vector sono nulli. In fase di test
$h(\x) = \sign\big(\sum_i\alpha_id_i\,k(\x,\x_i)+b\big)$, e si memorizzano solo i support vector (che ricorda il
K-NN).

\begin{figure}[ht]
\centering
\begin{tikzpicture}[x=1cm, y=0.9cm]
\foreach \i/\y in {1/1.5,2/0.5,3/-0.5,4/-1.5} \node[neur] (x\i) at (0,\y) {$x_\i$};
\foreach \s/\y in {1/1.9,2/0.65,3/-0.6,4/-1.85} \node[draw, rounded corners=1.5pt, font=\scriptsize, minimum height=6mm] (k\s) at (3,\y) {$k(\x,\x_\s)$};
\foreach \i in {1,2,3,4} \foreach \s in {1,2,3,4} \draw[grigio] (x\i.east) -- (k\s.west);
\node[circle, draw, thick, minimum size=9mm] (s) at (6,0) {$\Sigma$};
\foreach \s in {1,2,3,4} \draw[->] (k\s.east) -- (s) node[pos=0.5, font=\scriptsize, fill=white, inner sep=1pt] {$\alpha_\s d_\s$};
\node[font=\small] (b) at (6,1.6) {$b$}; \draw[->] (b) -- (s);
\draw[->, thick] (s) -- (7.6,0) node[right, font=\small] {$\sign(\cdot)$};
\node[font=\footnotesize, align=center] at (3,-2.75) {una ``unità'' per support vector};
\end{tikzpicture}
\caption{Vista ``architetturale'' di una SVM: lo strato nascosto ha tante unità quanti sono i support vector, e
ciascuna calcola il kernel tra l'input e il proprio support vector.}
\label{fig:svm-architettura}
\end{figure}

Il numero di ``unità nascoste'' di questa vista è pari al \textbf{numero di support vector}. Non sono però unità
di una rete neurale. In una rete le unità nascoste hanno \textbf{parametri adattivi}. Qui i ``centri'' sono i
support vector stessi, e la funzione è fissata dal kernel (\figrif{fig:svm-architettura}).

\subsubsection{Esempi di kernel}
\begin{itemize}
  \item \textbf{Polinomiale}: $k(\x,\x_i) = (\x\T\x_i+1)^p$, con $p$ scelto dall'utente.
  \item \textbf{RBF} (\ingl{Radial Basis Function}, kernel gaussiano):
        $k(\x,\x_i) = \exp\!\big(-\frac{1}{2\sigma^2}\norm{\x-\x_i}^2\big)$, con $\sigma^2$ scelto dall'utente. Si
        scrive anche con $\gamma = 1/(2\sigma^2)$. Con kernel molto ``stretti'' ($\sigma$ piccolo) la risposta su
        $\x_i$ è solo $d_i$, come nel 1-NN.
  \item \textbf{Perceptron a due strati}: $k(\x,\x_i) = \tanh(\beta_0\,\x\T\x_i+\beta_1)$, con $\beta_0>0$ e
        $\beta_1<0$.
\end{itemize}

\begin{attenzione}
Il kernel polinomiale e l'RBF sono \textbf{sempre} kernel validi (prodotti scalari). Per il perceptron a due
strati invece il teorema di Mercer vale solo per alcune scelte di $\beta_0$ e $\beta_1$. Un kernel RBF corrisponde
a uno spazio delle feature di \textbf{dimensione infinita}.
\end{attenzione}

\section{SVM per la regressione}

Si ha $d = f(\x)+v$, con $f$ sconosciuta e rumore $v$ indipendente da $\x$. Si stima $d$ con un'espansione
lineare di funzioni non lineari:
\[
h(\x) = \w\T\Phi(\x),\qquad \w = (b,w_1,\dots,w_{n_\Phi})\T,\quad \Phi(\x) = (1,\phi_1(\x),\dots,\phi_{n_\Phi}(\x))\T .
\]

\begin{attenzione}[Dove sta il bias]
In questa forma il bias è la componente $w_0=b$ di $\w$, con $\phi_0(\x)=1$. Il termine $\frac12\w\T\w$ del
problema qui sotto penalizza quindi anche $b$, e $b$ si ottiene insieme agli altri pesi. Se invece si tiene $b$
fuori da $\w$, come nella classificazione, va calcolato a parte dopo aver risolto il duale.
\end{attenzione}

\begin{definizione}[Loss $\eps$-insensitive]
\[
L_\eps(d,h(\x)) = \begin{cases}\abs{d-h(\x)}-\eps & \text{se }\abs{d-h(\x)}\ge\eps\\ 0 & \text{altrimenti.}\end{cases}
\]
\end{definizione}

Gli errori più piccoli di $\eps$ \textbf{non costano nulla}. Si forma un \textbf{tubo} di ampiezza $\eps$ attorno alla
predizione, e solo i punti fuori dal tubo contribuiscono alla loss (\figrif{fig:eps}). Di nuovo \textbf{solo i
support vector guidano la soluzione}. Nella classificazione erano i punti sul bordo o oltre il margine; qui sono i
punti sul bordo o fuori dal tubo.

\begin{figure}[ht]
\centering
\begin{tikzpicture}
\begin{axis}[stile, width=5.4cm, height=4cm, xmin=-2.2, xmax=2.2, ymin=0, ymax=1.6, xtick={-0.6,0.6},
  xticklabels={$-\eps$,$\eps$}, ytick=\empty, xlabel={$d-h(\x)$}, ylabel={$L_\eps$}]
\addplot[very thick, blu] coordinates {(-2.1,1.5) (-0.6,0) (0.6,0) (2.1,1.5)};
\end{axis}
\end{tikzpicture}\hspace{0.6cm}
\begin{tikzpicture}
\begin{axis}[stile, width=7cm, height=4.2cm, domain=0:6.3, samples=100, xmin=0, xmax=6.4, ymin=-1.9, ymax=1.9,
  xtick=\empty, ytick=\empty, xlabel={$x$}, ylabel={$d$}]
\addplot[name path=su, thin, densely dashed] {sin(deg(x))+0.45};
\addplot[name path=giu, thin, densely dashed] {sin(deg(x))-0.45};
\addplot[grigiochiaro] fill between[of=su and giu];
\addplot[thick, inchiostro] {sin(deg(x))};
\addplot[only marks, mark=*, mark size=1.5pt, blu] coordinates {(0.5,0.6) (1.2,0.8) (2.1,0.95) (3.0,0.0) (3.6,-0.3) (4.4,-1.1) (5.3,-0.75) (6.0,-0.4)};
\addplot[only marks, mark=*, mark size=1.5pt, rosso] coordinates {(1.6,1.75) (4.0,-1.6) (2.6,-0.2)};
\draw[<->, rosso] (axis cs:1.6,1.75) -- (axis cs:1.6,1.45) node[midway, right, font=\scriptsize] {$\xi$};
\draw[<->, rosso] (axis cs:4.0,-1.6) -- (axis cs:4.0,-1.207) node[midway, right, font=\scriptsize] {$\xi'$};
\draw[<->, rosso] (axis cs:2.6,-0.2) -- (axis cs:2.6,0.0655) node[midway, left, font=\scriptsize] {$\xi'$};
\end{axis}
\end{tikzpicture}
\caption{A sinistra la loss $\eps$-insensitive. A destra il tubo di ampiezza $\eps$ attorno alla funzione stimata:
solo i punti fuori dal tubo (rosso) hanno variabili slack non nulle.}
\label{fig:eps}
\end{figure}

Si introducono due insiemi di variabili slack non negative, $\xi_i$ e $\xi'_i$ (sopra e sotto il tubo), con
$-\xi'_i-\eps\le d_i-\w\T\Phi(\x_i)\le\eps+\xi_i$.

\begin{definizione}[Problema primale --- regressione]
Minimizzare
\[
\Psi(\w,\bxi,\bxi') = \frac12\,\w\T\w + C\sum_{i=1}^{l}(\xi_i+\xi'_i)
\]
soggetto a $d_i - \w\T\Phi(\x_i)\le\eps+\xi_i$, $\ \w\T\Phi(\x_i)-d_i\le\eps+\xi'_i$, $\ \xi_i\ge0$, $\ \xi'_i\ge0$.
\end{definizione}

\begin{definizione}[Problema duale --- regressione$^*$]
Trovare $\{\alpha_i\}$ e $\{\alpha'_i\}$ che massimizzano
\[
Q(\balpha,\balpha') = \sum_{i=1}^{l}d_i(\alpha_i-\alpha'_i) - \eps\sum_{i=1}^{l}(\alpha_i+\alpha'_i)
- \frac12\sum_{i,j=1}^{l}(\alpha_i-\alpha'_i)(\alpha_j-\alpha'_j)\,k(\x_i,\x_j)
\]
soggetto a $\sum_i(\alpha_i-\alpha'_i) = 0$ e $0\le\alpha_i,\alpha'_i\le C$.
\end{definizione}

Dalla soluzione, con $\gamma_i = \alpha_i-\alpha'_i$,
\[
\w = \sum_{i=1}^{l}\gamma_i\,\Phi(\x_i),\qquad h(\x) = \sum_{i=1}^{l}\gamma_i\,\Phi\T(\x_i)\Phi(\x) = \sum_{i=1}^{l}\gamma_i\,k(\x_i,\x),
\]
e i \textbf{support vector} corrispondono ai $\gamma_i$ non nulli. La procedura è la stessa della classificazione.
Si scelgono $C$, $\eps$ e il kernel; si calcola $K$; si risolve il duale; si calcola il bias, se è tenuto fuori da
$\w$. La funzione stimata è una combinazione lineare di prodotti scalari in uno spazio delle feature che non serve
conoscere.

\section{Le SVM nella pratica}

\begin{idea}[Pregi]
\begin{itemize}
  \item \textbf{Regolarizzazione incorporata} nel problema di ottimizzazione (il margine) e approssimazione della SRM:
        il bound sulla VC-dimension è minimizzato dalla soluzione, nel caso hard margin senza kernel.
  \item \textbf{Problema convesso}: l'addestramento trova sempre il \textbf{minimo globale}, a differenza delle reti
        neurali.
  \item Un modello lineare, semplice e compatto, con un ricco insieme di funzioni di decisione non lineari grazie alla
        \textbf{trasformazione implicita delle feature}: un trattamento efficiente della LBE. Un insieme ricco non
        implica una VC-dimension alta, che dipende da margine, parametri del kernel e $C$.
  \item Il kernel trick si estende ad altri modelli (\rif{sec:metodi-kernel}).
\end{itemize}
\end{idea}

\begin{attenzione}[Difetti]
\begin{itemize}
  \item Bisogna \textbf{scegliere il kernel e i suoi parametri}.
  \item È un algoritmo \textbf{batch}.
  \item La matrice kernel è $l\times l$. Problemi molto grandi (oltre circa 20\,000 esempi) erano difficili da
        risolvere esattamente con gli approcci standard. Oggi esistono molte soluzioni, anche basate sulla discesa
        del gradiente.
\end{itemize}
\end{attenzione}

\begin{figure}[ht]
\centering
\begin{tikzpicture}
\begin{groupplot}[group style={group size=2 by 1, horizontal sep=0.45cm}, scatola, width=6.4cm, height=5.8cm,
  xmin=-3, xmax=3.8, ymin=-2.6, ymax=3.4, xtick=\empty, ytick=\empty]
\nextgroupplot[title={\small kernel polinomiale (grado 4)}]
\addplot[only marks, mark=o, mark size=1.4pt, rosso!80!black, thin] table[x=x, y=y, restrict expr to domain={\thisrow{c}}{0:0}] {dati/htf_punti.dat};
\addplot[only marks, mark=*, mark size=1.2pt, blu] table[x=x, y=y, restrict expr to domain={\thisrow{c}}{1:1}] {dati/htf_punti.dat};
\addplot[only marks, mark=*, mark size=0.55pt, inchiostro] table {dati/svm_poly_sv.dat};
\addplot[thin, inchiostro, densely dashed] table {dati/svm_poly_marg.dat};
\addplot[thick, inchiostro] table {dati/svm_poly_conf.dat};
\addplot[thick, viola, densely dotted] table {dati/htf_bayes.dat};
\nextgroupplot[title={\small kernel RBF ($\gamma=1$)}]
\addplot[only marks, mark=o, mark size=1.4pt, rosso!80!black, thin] table[x=x, y=y, restrict expr to domain={\thisrow{c}}{0:0}] {dati/htf_punti.dat};
\addplot[only marks, mark=*, mark size=1.2pt, blu] table[x=x, y=y, restrict expr to domain={\thisrow{c}}{1:1}] {dati/htf_punti.dat};
\addplot[only marks, mark=*, mark size=0.55pt, inchiostro] table {dati/svm_rbf_sv.dat};
\addplot[thin, inchiostro, densely dashed] table {dati/svm_rbf_marg.dat};
\addplot[thick, inchiostro] table {dati/svm_rbf_conf.dat};
\addplot[thick, viola, densely dotted] table {dati/htf_bayes.dat};
\end{groupplot}
\end{tikzpicture}
\caption{SVM ($C=1$) sul problema a due classi: confine (nero), margini $g(\x)=\pm1$ (tratteggio), support vector
(punto nero al centro del simbolo) e confine di Bayes (viola). Polinomiale: errore di training $\valsvmpolyTrain$,
di test $\valsvmpolyTest$, support vector $\valsvmpolySV\%$ dei dati. RBF: training $\valsvmrbfTrain$, test
$\valsvmrbfTest$, support vector $\valsvmrbfSV\%$. Errore di Bayes $\valbayesErr$.}
\label{fig:svm-pratica}
\end{figure}

Sul problema a due classi (\figrif{fig:svm-pratica}) il kernel radiale dà il risultato migliore. È quello che ci
si può aspettare, visto che i dati provengono da miscele di gaussiane. Si comporta in modo simile alla rete
regolarizzata della \figrif{fig:weight-decay}. I support vector sono i punti sul margine (slack nulla) o dalla
parte sbagliata del margine (slack positiva).

\begin{attenzione}[I support vector possono essere tanti]
La teoria immagina pochi support vector, ma \textbf{non è sempre vero}: qui sono circa la metà dei dati. La sparsità
della soluzione non è garantita.
\end{attenzione}

\subsection{Il ruolo degli iperparametri}

Non esiste una teoria che garantisca che una data famiglia di SVM abbia alta accuratezza su un dato problema. Le
proprietà del caso hard margin \textbf{non si estendono} direttamente al soft margin e ai kernel: $C$ e il kernel
possono portare a una VC-dimension anche \textbf{infinita}. Si prenda una SVM con kernel gaussiano di larghezza
$\sigma$ \textbf{abbastanza piccola}. Classifica correttamente un numero arbitrario di punti, come il 1-NN, perché
contribuisce solo il support vector più vicino. Ha quindi VC-dimension infinita, a meno di regolarizzare. Con una
larghezza \textbf{grande} invece contribuiscono tutti i support vector. Si ottiene una sorta di media globale, con
VC-dimension bassa. \uline{Controllando la larghezza del kernel si controlla la VC-dimension.}

\begin{figure}[ht]
\centering
\begin{tikzpicture}
\begin{axis}[stile, width=8.4cm, height=5cm, xmode=log, xmin=0.01, xmax=1000, ymin=0.2, ymax=0.42, xlabel={$C$},
  ylabel={errore di test}, ytick={0.2,0.25,0.3,0.35,0.4}, yticklabels={$0{,}20$,$0{,}25$,$0{,}30$,$0{,}35$,$0{,}40$},
  legend pos=outer north east]
\addplot[thick, rosso, mark=o, mark size=1.2pt] table[x=C, y=g5] {dati/svm_C_gamma.dat}; \addlegendentry{$\gamma=5$}
\addplot[thick, arancio, mark=o, mark size=1.2pt] table[x=C, y=g1] {dati/svm_C_gamma.dat}; \addlegendentry{$\gamma=1$}
\addplot[thick, verde, mark=o, mark size=1.2pt] table[x=C, y=g05] {dati/svm_C_gamma.dat}; \addlegendentry{$\gamma=0{,}5$}
\addplot[thick, blu, mark=o, mark size=1.2pt] table[x=C, y=g01] {dati/svm_C_gamma.dat}; \addlegendentry{$\gamma=0{,}1$}
\end{axis}
\end{tikzpicture}
\caption{Errore di test di una SVM con kernel RBF in funzione di $C$ per diversi valori di $\gamma=1/(2\sigma^2)$, sul
problema a due classi: il valore migliore di $C$ cambia con $\gamma$.}
\label{fig:svm-C-gamma}
\end{figure}

Con kernel stretti ($\gamma$ grande) conviene un $C$ basso, cioè più regolarizzazione. Con $\gamma$ piccoli l'ottimo
si sposta verso $C$ più alti (\figrif{fig:svm-C-gamma}). \textbf{Il miglior valore di $C$ dipende dagli
iperparametri del kernel}: vanno selezionati \textbf{insieme} (grid search). Serve una validazione rigorosa
(validation e test set, oppure double CV), perché la loro regolazione può introdurre overfitting.

\begin{attenzione}[Luoghi comuni da sfatare]
\begin{enumerate}
  \item \emph{``La SVM si può usare con i valori di default.''} No: $C$ e i parametri del kernel sono critici.
  \item \emph{``Scelgo la SVM perché non va in overfitting.''} No: con soft margin e kernel la VC-dimension può essere
        altissima. L'overfitting dipende da $C$ e dal kernel.
  \item \emph{``La SVM risolve la maledizione della dimensionalità.''} No: gestisce bene l'alta dimensione dello
        \textbf{spazio delle feature}, non dello \textbf{spazio di input}.
  \item \emph{``La SVM trova da sola il numero di unità.''} Il numero di support vector è determinato dalla
        soluzione. Ma dipende da $C$ e dal kernel, che vanno scelti, e non è detto che sia piccolo.
\end{enumerate}
\end{attenzione}

\section{Metodi kernel}
\label{sec:metodi-kernel}

Il kernel trick si usa anche senza SVM: è la \termine{kernelizzazione} di approcci tradizionali. Si supponga che in
un modello compaia un \textbf{prodotto scalare} o una \textbf{misura di similarità}. Lo si può sostituire con un
kernel. Il modello opera allora in un nuovo spazio implicito ad alta dimensione, semplicemente specificando il
kernel. È un'espansione in basi fissa e ad alta dimensione, ma efficiente. La scelta delle basi è
\textbf{modulare}: si cambia il kernel senza cambiare la macchina di apprendimento. E il quadro teorico è unificato.

\begin{proprieta}{Kernel e distanza}
Un kernel induce una distanza nello spazio delle feature:
\[
d_k(\vv s,\vv t)^2 = \norm{\Phi(\vv s)-\Phi(\vv t)}^2 = k(\vv s,\vv s) - 2k(\vv s,\vv t) + k(\vv t,\vv t).
\]
\end{proprieta}
\begin{proof}
Con lo sviluppo \eqref{eq:quadrato-norma}:
$\norm{\Phi(\vv s)-\Phi(\vv t)}^2 = \Phi(\vv s)\T\Phi(\vv s) - 2\,\Phi(\vv s)\T\Phi(\vv t) + \Phi(\vv t)\T\Phi(\vv t)$, e
ogni prodotto scalare è un valore del kernel.
\end{proof}

Il kernel è quindi una misura di \textbf{similarità}: oggetti simili (distanza piccola) danno un valore alto del
kernel.

\begin{idea}[L'idea dei kernel]
Ci si concentra sulla rappresentazione dei \textbf{confronti tra coppie} di oggetti, invece che sulla
rappresentazione di un singolo oggetto. In alcuni domini è più facile confrontare due oggetti che definire uno
spazio di feature in cui rappresentarne uno. La conoscenza di dominio sta tutta nel kernel, \textbf{separata dal
risolutore}.
\end{idea}

Un kernel è buono quando raggruppa nello spazio delle feature gli oggetti della stessa classe, cioè quando riflette
la similarità rilevante per il task. È cattivo quando la similarità che induce non ha a che fare con il task, ad
esempio se tutti gli oggetti risultano ugualmente simili o ugualmente diversi. Una SVM è caratterizzata in gran
parte dalla scelta del kernel. Nelle reti neurali invece le funzioni di base si \textbf{adattano} ai dati. La
scelta del kernel migliore per un problema è un tema di ricerca (kernel adattivi). La progettazione di kernel si
estende a oggetti non vettoriali, come stringhe, alberi e grafi (\rif{sec:kernel-strutture}).

\begin{puntichiave}
\begin{itemize}
  \item Vincoli canonici $d_i(\w\T\x_i+b)\ge1$, margine $\rho = 2/\norm{\w}$: massimizzare il margine = minimizzare
        $\frac12\w\T\w$.
  \item Duale: $\max\sum\alpha_i - \frac12\sum\alpha_i\alpha_jd_id_j\,\x_i\T\x_j$; $\w = \sum\alpha_id_i\x_i$ dipende solo
        dai support vector ($\alpha_i>0$).
  \item Teorema di Vapnik: VC $\le\min(\lceil D^2/\rho^2\rceil,n)+1$; margine grande, VC piccola.
  \item Soft margin: slack $\xi_i$, penalità $C\sum\xi_i$, duale con $0\le\alpha_i\le C$.
  \item Kernel $k(\x,\y) = \Phi\T(\x)\Phi(\y)$ (Mercer: matrice semidefinita positiva); RBF di dimensione infinita.
  \item Regressione: loss $\eps$-insensitive, $h(\x) = \sum\gamma_ik(\x_i,\x)$.
  \item $C$ e parametri del kernel vanno scelti insieme; la VC-dimension può essere infinita.
\end{itemize}
\end{puntichiave}

\begin{domande}
\begin{enumerate}
  \item Definire il classificatore a margine massimo e ricavare $\rho = 2/\norm{\w}$.
  \item Come il margine massimo diventa un problema di programmazione quadratica? Cosa sono i support vector?
  \item Perché la SVM approssima la SRM?
  \item Come si gestiscono dati non separabili? Ruolo di $C$.
  \item Come si ottengono confini non lineari? Cos'è il kernel trick? Perché permette di lavorare con espansioni in
        basi di dimensione altissima?
  \item Come funziona la SVM per la regressione?
  \item Perché la VC-dimension di una SVM con kernel RBF può essere infinita? Perché $C$ e i parametri del kernel vanno
        selezionati insieme? La SVM risolve la maledizione della dimensionalità?
  \item Cosa sono i metodi kernel? Che legame c'è tra kernel e similarità?
\end{enumerate}
\end{domande}
